\medskip

Integrace interaktivních AI metod do výuky má dvě vzájemně se doplňující roviny: pedagogické cíle (co chceme u studentů podpořit) a praktické nástroje/scénáře, které to umožní. Níže nejprve shrnujeme hlavní očekávané přínosy (stručně jako doporučené cíle) a poté uvádíme konkrétní nízkoprahové příklady, které lze v předmětu otestovat.

\smallskip
\textbf{Doporučené pedagogické cíle (shrnutí, general background knowledge)}  
\begin{itemize}
  \item Zvýšení angažovanosti studentů skrze interaktivní, praktické aktivity a simulace. (general background knowledge)
  \item Zrychlení a škálování formativní zpětné vazby pomocí AI nástrojů, které generují komentáře nebo diagnostiku chyb. (general background knowledge)
  \item Podpora diferenciace učení — adaptivní úlohy a personalizované nápovědy pro různé úrovně studentů. (general background knowledge)
  \item Důraz na proces (dokumentace postupu, milníky) místo pouze konečného výsledku; integrace reflexe o použití AI do hodnocení. (general background knowledge)
\end{itemize}

\smallskip
\textbf{Nízkoprahové praktické příklady — co zkusit v prvním pilotu}  

\begin{enumerate}
  \item Minecraft Education — názorné demonstrace principů AI a „Hour of Code“ aktivity.  
    Minecraft Education obsahuje připravené světy zaměřené i na principy fungování AI a zdarma dostupné hodiny kódu, které lze použít pro rychlou ukázku a hands‑on aktivity ve výuce \cite{kb_ai_101_tipu_pdf_04e4a115_page_15_2}. Praktický krok: stáhněte Minecraft Education a vyzkoušejte se studenty jednu ze zdarma dostupných hodin kódu; po aktivitě zařaďte krátkou reflexi, kde studenti popíšou, co agent „pozoroval“ a jak se chování změnilo \cite{kb_ai_101_tipu_pdf_04e4a115_page_15_2}.
  \item Procvičování s diagnostikou chyb — „Pokrok v matematice“.  
    Funkce „Pokrok v matematice“ umožňuje učiteli vygenerovat sadu příkladů, vyžadovat nahrání postupu řešení a získat přehled o oblastech, kde studenti chybují; nástroj také předopraví a doporučí oblast s chybou, čímž usnadňuje cílenou intervenci \cite{kb_ai_101_tipu_pdf_04e4a115_page_8_1}.
  \item Rychlá tvorba elektronických testů — Microsoft 365 Copilot + Forms.  
    Kombinace Microsoft 365 Copilot a Forms umí automaticky generovat písemky/kvízy (např. mix ABC, pravda/nepravda, otevřené otázky) včetně označení správných odpovědí a krátkých vysvětlení — užitečné pro rychlé, škálovatelné hodnocení v kontrolovaných podmínkách \cite{kb_ai_101_tipu_pdf_04e4a115_page_36_1}.
\end{enumerate}

\smallskip
\textbf{Krátký plán prvního interaktivního sezení (praktický checklist)}  
\begin{enumerate}
  \item Vyberte jeden jasný cíl sezení (např. „ukázat princip učení agenta“ nebo „rychlá diagnostika chyb v řešení rovnic“). (general background knowledge)
  \item Zvolte nástroj podle cíle: Minecraft Education pro demonstrace (viz výše) \cite{kb_ai_101_tipu_pdf_04e4a115_page_15_2}, „Pokrok v matematice“ pro procvičování s požadavkem na postupy řešení \cite{kb_ai_101_tipu_pdf_04e4a115_page_8_1}, nebo Copilot+Forms pro generování kvízu \cite{kb_ai_101_tipu_pdf_04e4a115_page_36_1}.
  \item Připravte krátký instrukční list pro studenty (co se očekává, jak deklarovat použití AI) a jednu reflexní otázku po aktivitě. (general background knowledge)
  \item Po pilotu zaznamenejte data a zpětnou vazbu: co fungovalo, jaké technické/etické problémy přišly, návrh další iterace. (general background knowledge)
\end{enumerate}

\smallskip
Poznámka k rizikům a řízení kvality: konkrétní rizika (ochrana osobních údajů, akademická integrita, nekonzistentní nebo nesprávné výstupy AI) je třeba před pilotem zhodnotit a popsat v krátkém provozním plánu; tato doporučení jsou rozvedena v kapitolách věnovaných rizikům, GDPR a akademické integritě (viz příručka). (general background knowledge)